% Part: normal-modal-logic
% Chapter: filtrations
% Section: finite

\documentclass[../../../include/open-logic-section]{subfiles}

\begin{document}

\olfileid{nml}{fil}{fin}

\olsection{निस्यंदने सांत असतात}

उप-!!{formula}s खाली बंद असलेल्या कोणत्याही संच~$\Gamma$
साठी आपण निस्यंदनांची व्याख्या केली. केवळ त्या
व्याख्येवरून निस्यंदने सांत असतील याची खात्री मिळत नाही.
खरे तर~$\Gamma$ अनंत असेल (उदाहरणार्थ, सर्व
!!{formula}s चा संच असेल), तर निस्यंदनही अनंत असू शकते.
मात्र~$\Gamma$ सांत असेल (उदाहरणार्थ, दिलेल्या
!!{formula}~$!A$ च्या उप-!!{formula}s चा संच असेल),
तर~$\Gamma$ मधून मिळणारे प्रत्येक निस्यंदनही सांत असते.

\begin{prop}\ollabel{prop:filt-are-finite}
  $\Gamma$ सांत असेल, तर प्रतिमान~$\mModel{M}$ चे
  $\Gamma$ मधून मिळणारे कोणतेही निस्यंदन~$\mModel{M^*}$
  सुद्धा सांत असते.
\end{prop}

\begin{proof}
  $W^*$ चा आकार म्हणजे तुल्यता संबंध~$\equiv$ अंतर्गत
  वेगवेगळ्या वर्ग~$[w]$ ची संख्या. अशा एकाच वर्गातील
  कोणतीही दोन जगे $u$, $v$, म्हणजे $u$ आणि $v$
  अशी जगे की $u \equiv v$, ही $\Gamma$ मधील सर्व !!{formula}s~$!A$
  यांच्या सत्यतेबाबत एकमत असतात: $!A \in \Gamma$ असेल,
  तर~$!A$ हे $u$ आणि~$v$ या दोन्ही जगांवर सत्य असते
  किंवा दोन्हीवर असत्य असते. म्हणून प्रत्येक वर्ग~$[w]$
  हा~$\Gamma$ च्या एका उपसंचाशी जुळतो:~$[w]$ मधील
  जगांवर~$!A$ सत्य असेल अशा सर्व $!A \in \Gamma$ च्या संचाशी.
  दोन वेगवेगळे वर्ग $[u]$ आणि $[v]$ हे $\Gamma$ च्या
  एकाच उपसंचाशी जुळू शकत नाहीत. कारण~$u$ येथे सत्य
  असलेल्या !!{formula}s चा संच आणि~$v$ येथे सत्य असलेल्या
  !!{formula}s चा संच एकच असेल, तर~$u$ आणि~$v$ यांच्यात
  $\Gamma$ मधील सर्व सूत्रांच्या सत्यतेबाबत एकमत असते;
  म्हणजे $u \equiv v$. पण मग $[u] = [v]$.
  म्हणून $W^*$ पासून~$\Pow{\Gamma}$ मध्ये !!a{injective}
  फल आहे आणि त्यामुळे $\card{W^*} \le \card{\Pow{\Gamma}}$.
  म्हणून~$\Gamma$ मध्ये $n$ वाक्ये असतील, तर
  $W^*$ ची गणनसंख्या~$2^n$ पेक्षा अधिक नाही.
\end{proof}

\end{document}
